\documentclass[11pt,letterpaper]{article}

% ============================================================
% Page layout and typography
% ============================================================
\usepackage[
  top=0.68in,
  bottom=0.68in,
  left=0.80in,
  right=0.80in
]{geometry}
\usepackage[T1]{fontenc}
\usepackage{lmodern}
\usepackage{microtype}
\usepackage{amsmath,amssymb,mathtools}
\usepackage{booktabs,array,tabularx}
\usepackage{xcolor}
\usepackage{listings}
\usepackage{enumitem}
\usepackage{titlesec}
\usepackage{fancyhdr}
\usepackage[most]{tcolorbox}
\usepackage{xurl}
\usepackage{hyperref}

% ============================================================
% Colors
% ============================================================
\definecolor{wvuBlue}{RGB}{0,40,85}
\definecolor{wvuGold}{RGB}{234,170,0}
\definecolor{codebg}{RGB}{247,248,250}
\definecolor{codeframe}{RGB}{205,210,216}
\definecolor{codecomment}{RGB}{55,105,65}
\definecolor{codestring}{RGB}{145,55,45}
\definecolor{notebg}{RGB}{247,250,253}

% ============================================================
% Hyperlinks
% ============================================================
\hypersetup{
  colorlinks=true,
  linkcolor=wvuBlue,
  urlcolor=wvuBlue,
  citecolor=wvuBlue,
  pdfauthor={MAE 623: Heat Conduction},
  pdftitle={Transient Heat Conduction in Python}
}
\urlstyle{same}

% ============================================================
% Paragraphs, lists, and headings
% ============================================================
\setlength{\parindent}{0pt}
\setlength{\parskip}{4pt}
\setlength{\emergencystretch}{3em}
\raggedbottom

\setlist[itemize]{leftmargin=0.28in,itemsep=1.5pt,topsep=2.5pt}
\setlist[enumerate]{leftmargin=0.30in,itemsep=2pt,topsep=2.5pt}

\titleformat{\section}
  {\Large\bfseries\color{wvuBlue}}
  {\thesection.}{0.55em}{}
\titleformat{name=\section,numberless}
  {\Large\bfseries\color{wvuBlue}}
  {}{0pt}{}
\titleformat{\subsection}
  {\large\bfseries\color{wvuBlue}}
  {}{0pt}{}
\titlespacing*{\section}{0pt}{12pt}{5pt}
\titlespacing*{\subsection}{0pt}{9pt}{3pt}

% ============================================================
% Header and footer
% ============================================================
\pagestyle{fancy}
\fancyhf{}
\fancyhead[L]{\small MAE 623: Heat Conduction}
\fancyhead[R]{\small Transient Heat Conduction in Python}
\fancyfoot[C]{\small \thepage}
\renewcommand{\headrulewidth}{0.4pt}
\renewcommand{\footrulewidth}{0pt}
\setlength{\headheight}{14pt}

% ============================================================
% Tables
% ============================================================
\newcolumntype{P}[1]{>{\raggedright\arraybackslash}p{#1}}
\newcolumntype{Y}{>{\raggedright\arraybackslash}X}
\renewcommand{\arraystretch}{1.14}

% ============================================================
% Code listings
% ============================================================
\lstdefinestyle{python}{
  language=Python,
  basicstyle=\ttfamily\small,
  keywordstyle=\color{wvuBlue},
  commentstyle=\color{codecomment},
  stringstyle=\color{codestring},
  backgroundcolor=\color{codebg},
  frame=single,
  framerule=0.4pt,
  rulecolor=\color{codeframe},
  columns=fullflexible,
  keepspaces=true,
  showstringspaces=false,
  breaklines=true,
  breakatwhitespace=false,
  tabsize=4,
  xleftmargin=0.01\linewidth,
  xrightmargin=0.01\linewidth,
  framexleftmargin=5pt,
  framexrightmargin=5pt,
  aboveskip=5pt,
  belowskip=5pt
}

\lstdefinestyle{terminal}{
  basicstyle=\ttfamily\small,
  backgroundcolor=\color{codebg},
  frame=single,
  framerule=0.4pt,
  rulecolor=\color{codeframe},
  columns=fullflexible,
  keepspaces=true,
  showstringspaces=false,
  breaklines=true,
  xleftmargin=0.01\linewidth,
  xrightmargin=0.01\linewidth,
  framexleftmargin=5pt,
  framexrightmargin=5pt,
  aboveskip=5pt,
  belowskip=5pt
}

\lstset{style=python}

% ============================================================
% Reusable notation and callout boxes
% ============================================================
\newcommand{\Fo}{\mathrm{Fo}}
\newcommand{\Bi}{\mathrm{Bi}}
\newcommand{\code}[1]{\texttt{\detokenize{#1}}}

\newtcolorbox{guidebox}[1][]{
  enhanced,
  breakable,
  colback=notebg,
  colframe=wvuBlue!55,
  boxrule=0.6pt,
  arc=1.5mm,
  left=8pt,
  right=8pt,
  top=7pt,
  bottom=7pt,
  before skip=7pt,
  after skip=8pt,
  #1
}

\begin{document}

% ============================================================
% Title block
% ============================================================
\begin{center}
  {\LARGE\bfseries\color{wvuBlue} Transient Heat Conduction in Python}\par
  \vspace{4pt}
  {\large A Student Guide to Explicit and Implicit Methods}\par
  \vspace{1pt}
  {\large and Changing Boundary Conditions}\par
  \vspace{5pt}
  {\normalsize Companion to the two one-dimensional example scripts}\par
  \vspace{8pt}
  {\color{wvuGold}\rule{\textwidth}{1.1pt}}
\end{center}
\vspace{2pt}
\section*{What you will learn}

This guide connects the heat equation to the calculations in
\code{transient_1d_explicit.py} and \code{transient_1d_implicit.py}.
No previous Python programming experience is assumed.

\begin{guidebox}
By the end of the guide, you should be able to:
\begin{itemize}
  \item interpret the arrays, indexing, and time-stepping loops used in both scripts;
  \item explain why the explicit method updates temperatures directly while the implicit method solves a system of simultaneous equations; and
  \item replace a fixed-temperature boundary with a specified heat flux, insulation, or convection boundary condition.
\end{itemize}
\end{guidebox}

\begin{guidebox}
\textbf{Use this guide with copies of the scripts.}
The programs currently hold the left end at $100\,^{\circ}\mathrm C$ and the right end at
$20\,^{\circ}\mathrm C$. The worked boundary changes replace the right condition while retaining the fixed left temperature. These boundary recipes are alternatives, so apply one recipe at a time. The original scripts are not modified by this document.
\end{guidebox}


\section{Run a script and find the inputs}
Install Python 3 from \url{https://www.python.org/downloads/} if needed.
On Windows, select the installer option to add Python to PATH.
Save both Python files in a folder. Open that folder in File Explorer,
type \code{cmd} in its address bar, and press Enter. In that command
window, install the packages once and run either script:
\begin{lstlisting}[style=terminal]
python -m pip install numpy matplotlib
python transient_1d_explicit.py
python transient_1d_implicit.py
\end{lstlisting}
Run the commands one at a time when comparing timings. Each program prints
its solver time and saves a contour PNG beside the script; no plot window
opens. The batch launcher runs both programs concurrently, which can affect
timings. Edit code in a plain-text editor, save it, and rerun the command.


The editable inputs are near the top of each script. Use the same physical
inputs and grid in both files for a direct comparison. The current settings
are $k=45\ \mathrm{W/(m\,K)}$, $\rho=7800\ \mathrm{kg/m^3}$,
$c_p=500\ \mathrm{J/(kg\,K)}$, $L=0.10\ \mathrm m$, 201 nodes,
$\Delta t=0.005\ \mathrm s$, and a final time of $300\ \mathrm s$.
For quicker practice with boundary changes, use 51 nodes and
\code{dt = 0.1} in both copies; this also satisfies the original explicit
stability limit.


\section{The physical equation and the grid}
For a uniform rod with constant properties, no internal heat generation,
and no heat loss through its sides,
\begin{equation}
 \rho c_p\frac{\partial T}{\partial t}
 =k\frac{\partial^2T}{\partial x^2},\qquad
 \frac{\partial T}{\partial t}
 =\alpha\frac{\partial^2T}{\partial x^2},\qquad
 \alpha=\frac{k}{\rho c_p}.
\end{equation}
The left side describes heat storage; the right side describes the net
effect of conduction. Thermal diffusivity $\alpha$ has units of
$\mathrm{m^2/s}$. Larger diffusivity makes temperature differences spread
through the rod more quickly.


Let $N=\code{number_of_nodes}-1$. Nodes are numbered $i=0,1,\ldots,N$,
so there are $N+1$ nodes but only $N$ intervals:
\[
 x_i=i\Delta x,\quad \Delta x=L/N,\qquad
 t_n=n\Delta t,\qquad T_i^n\approx T(x_i,t_n).
\]
The superscript $n$ labels a time; it is \emph{not} a power. The subscript
$i$ labels a position. Interior nodes have $1\le i\le N-1$.
The dimensionless Fourier number is
\[
 \Fo=\frac{\alpha\Delta t}{\Delta x^2}.
\]
With the current inputs, $\alpha\approx1.15385\times10^{-5}\ \mathrm{m^2/s}$,
$\Delta x=0.0005\ \mathrm m$, and $\Fo\approx0.230769$.


\subsection*{How the symbols become Python}
\begin{lstlisting}
alpha = k / (rho * cp)          # * multiplies; / divides
dx = L / (number_of_nodes - 1)
Fo = alpha * dt / dx**2         # **2 means squared
x = np.linspace(0.0, L, number_of_nodes)
time = np.arange(number_of_steps + 1) * dt
T_history = np.zeros((number_of_steps + 1, number_of_nodes))
\end{lstlisting}
\code{np} is the short name given to NumPy. An array is an organized
collection of numbers. Here, each row of \code{T_history} is a saved time
and each column is a position. Thus $T_i^n$ is \code{T_history[n, i]}.
There are 60,001 rows and 201 columns in the current full history, including
time zero. The temperature array alone occupies about 96 MB; plotting
uses additional memory.


\begin{center}
\small
\begin{tabularx}{0.95\linewidth}{@{}P{0.33\linewidth}Y@{}}
\toprule
\textbf{Python expression} & \textbf{Meaning} \\
\midrule
\code{T_history[0, :]} & Initial temperatures at all positions \\
\code{T_history[n, :]} & All temperatures at saved time $n$ \\
\code{T_history[:, 0]} & Left-end temperature at every time \\
\code{T_history[:, -1]} & Right-end temperature at every time \\
\code{T_history[n, -2]} & Node next to the right end at time $n$ \\
\bottomrule
\end{tabularx}
\end{center}
Python starts counting at zero. A colon selects all entries along an
array direction. Negative indices count backward: $-1$ is last, $-2$ is
next to last. Do not confuse Python index \code{-1} with a ghost node.


\section{Explicit FTCS: calculate from the old temperatures}
Approximate the spatial curvature by the centered difference
\[
 \left.\frac{\partial^2T}{\partial x^2}\right|_{i,n}
 \approx \frac{T_{i-1}^n-2T_i^n+T_{i+1}^n}{\Delta x^2}.
\]
Use a forward difference for the time derivative and evaluate the
curvature at the old time $n$:
\[
 \frac{T_i^{n+1}-T_i^n}{\Delta t}
 =\alpha\frac{T_{i-1}^n-2T_i^n+T_{i+1}^n}{\Delta x^2}.
\]
Rearranging gives the forward-time, centered-space (FTCS) update:
\begin{equation}\label{eq:explicit}
 \boxed{T_i^{n+1}=T_i^n+\Fo(T_{i-1}^n-2T_i^n+T_{i+1}^n).}
\end{equation}
Every quantity on the right is known. For example, with $\Fo=0.25$
and old left, center, and right temperatures $80,40,20\,^{\circ}\mathrm C$,
the new center temperature is $40+0.25(80-80+20)=45\,^{\circ}\mathrm C$.


\begin{lstlisting}
for n in range(number_of_steps):
    for i in range(1, number_of_nodes - 1):
        T_history[n + 1, i] = T_history[n, i] + Fo * (
            T_history[n, i - 1]
            - 2.0 * T_history[n, i]
            + T_history[n, i + 1]
        )
\end{lstlisting}
\code{for} repeats the indented instructions. \code{range(a, b)} produces
integers starting at \code{a} and stopping \emph{before} \code{b}.
The outer loop advances time. The inner loop updates interior nodes only.
Indentation is part of Python's syntax: use four spaces per level.
The right side reads row \code{n}; the left side writes row \code{n + 1}.
Changing row \code{n} in place would mix old and new temperatures and
would no longer be this method.


\subsection*{Why the time step is limited}
Equation~\eqref{eq:explicit} can be written
\[
 T_i^{n+1}=\Fo T_{i-1}^n+(1-2\Fo)T_i^n+\Fo T_{i+1}^n.
\]
For $0\le\Fo\le1/2$, the weights are nonnegative and sum to one:
the update is a weighted average. The standard one-dimensional FTCS
scheme with fixed end temperatures requires
\[
 \Fo\le\tfrac12,\qquad \Delta t\le\frac{\Delta x^2}{2\alpha}.
\]
The scripts stop if this condition is violated. Halving $\Delta x$ requires
reducing $\Delta t$ by a factor of four to keep the same Fourier number.
Refining the grid therefore increases both the nodes per step and the
number of steps. A convection boundary introduces an additional restriction
described later.


\section{Implicit Backward Euler: solve for the new temperatures}
Backward Euler evaluates spatial curvature at the \emph{new} time:
\[
 \frac{T_i^{n+1}-T_i^n}{\Delta t}
 =\alpha\frac{T_{i-1}^{n+1}-2T_i^{n+1}+T_{i+1}^{n+1}}{\Delta x^2}.
\]
Collecting the new temperatures on the left gives
\begin{equation}\label{eq:implicit}
 \boxed{-\Fo T_{i-1}^{n+1}+(1+2\Fo)T_i^{n+1}
        -\Fo T_{i+1}^{n+1}=T_i^n.}
\end{equation}
The new neighbors are unknown too. Solve all these equations together as
$A\mathbf T^{n+1}=\mathbf b$. For a five-node illustration, including the
two fixed endpoints, let $d=1+2\Fo$:
\[
\begin{bmatrix}
1&0&0&0&0\\
-\Fo&d&-\Fo&0&0\\
0&-\Fo&d&-\Fo&0\\
0&0&-\Fo&d&-\Fo\\
0&0&0&0&1
\end{bmatrix}
\begin{bmatrix}T_0^{n+1}\\T_1^{n+1}\\T_2^{n+1}\\T_3^{n+1}\\T_4^{n+1}\end{bmatrix}
=\begin{bmatrix}T_{\mathrm{left}}\\T_1^n\\T_2^n\\T_3^n\\T_{\mathrm{right}}\end{bmatrix}.
\]
The matrix is tridiagonal: only its main diagonal and the two neighboring
diagonals can contain nonzero entries. The example stores it as an ordinary
NumPy array for readability.
\begin{lstlisting}
A = np.zeros((number_of_nodes, number_of_nodes))
for i in range(1, number_of_nodes - 1):
    A[i, i - 1] = -Fo
    A[i, i] = 1.0 + 2.0 * Fo
    A[i, i + 1] = -Fo
A[0, 0] = 1.0
A[-1, -1] = 1.0


for n in range(number_of_steps):
    b = T_history[n, :].copy()
    b[0] = T_left
    b[-1] = T_right
    T_history[n + 1, :] = np.linalg.solve(A, b)
\end{lstlisting}
\code{A[i, j]} is the coefficient of unknown $j$ in equation $i$.
\code{copy()} makes an independent array so changing \code{b} does not
overwrite the old temperature row. \code{np.linalg.solve(A, b)} returns
the entire new row. There is no need to calculate a matrix inverse.


\section{Compare methods and recognize fixed boundaries}
\begin{center}
\small
\begin{tabularx}{\linewidth}{@{}P{0.22\linewidth}Y Y@{}}
\toprule
\textbf{Feature} & \textbf{Explicit FTCS} & \textbf{Implicit Backward Euler} \\
\midrule
Neighbors used & Old time $n$ & New time $n+1$ \\
Work each step & Direct temperature updates & Simultaneous linear solve \\
Stability for the original problem & Requires $\Fo\le0.5$ & Stable for any positive $\Delta t$ \\
Interior accuracy for smooth solutions & First order in time; second order in space & First order in time; second order in space \\
\bottomrule
\end{tabularx}
\end{center}
Unconditional stability does not mean unconditional accuracy. A large
implicit time step can smooth away a rapid transient. Both methods should
approach the same physical solution as space and time steps are reduced.
The sudden initial change at the left end is not smooth, so early-time
errors can be larger than later-time errors.


The printed time measures solver setup and time advancement, excluding
imports, shared grid/history allocation, and plotting. The explicit code
uses Python loops; the implicit code calls a dense NumPy solver at every
step. These timings compare the supplied implementations. A specialized
tridiagonal algorithm would change the implicit cost substantially.


\subsection*{Dirichlet means the temperature itself is prescribed}
The original boundary equations are
\[
 T_0^n=T_{\mathrm{left}},\qquad T_N^n=T_{\mathrm{right}}.
\]
In both scripts, the initialization block is
\begin{lstlisting}
T_history[0, :] = T_initial
T_history[:, 0] = T_left
T_history[:, -1] = T_right
\end{lstlisting}
The last two lines fill entire columns. Thus the endpoint values override
the uniform initial temperature at time zero. Physically, the left surface
is suddenly raised to $100\,^{\circ}\mathrm C$, while the interior starts
at $20\,^{\circ}\mathrm C$.


The explicit loop leaves these columns alone. The implicit code also
enforces them with identity matrix rows and the right-hand-side values
\code{b[0]} and \code{b[-1]}. To change a fixed right boundary, remove
\emph{all} of its fixed-temperature enforcement, not just one line.


\subsection*{A safe editing workflow}
Make a copy of the appropriate file and choose one right-boundary recipe
on the following pages. Replace the identified statements instead of
appending a second solver loop. Keep the grid, interior equations, timing
statements, and plotting section. A line beginning with \code{#} is a
comment; Python ignores it. Pasting an equation into a comment explains
it but does not change the calculation.


\section{Boundary physics: flux, insulation, and convection}
Take $x$ positive from left to right. At the right end, positive outward
heat flux removes heat from the rod:
\[
 -k\left.\frac{\partial T}{\partial x}\right|_L=q_{\mathrm{out}}.
\]
A negative $q_{\mathrm{out}}$ adds heat there. Insulation is the special
case $q_{\mathrm{out}}=0$. A specified flux determines a temperature
\emph{gradient}; it does not fix the surface temperature.


For convection at the right end,
\[
 q_{\mathrm{out}}=h(T_N-T_\infty),\qquad
 \Bi=\frac{h\Delta x}{k}.
\]
Here $h\ge0$ is the convection coefficient in $\mathrm{W/(m^2\,K)}$,
and $T_\infty$ is the surrounding fluid temperature. This $\Bi$ is a
\emph{mesh} Biot number based on $\Delta x$. A surface hotter than the
fluid loses heat; a colder surface gains heat.


\subsection*{Where the factor of two comes from}
Introduce a fictitious node $N+1$ just outside the right surface. A centered
approximation to the flux gives
\[
 -k\frac{T_{N+1}-T_{N-1}}{2\Delta x}=q_{\mathrm{out}}
 \quad\Longrightarrow\quad
 T_{N+1}=T_{N-1}-\frac{2\Delta x}{k}q_{\mathrm{out}}.
\]
Substituting this into the surface heat equation gives
\begin{equation}\label{eq:surface}
 \frac{dT_N}{dt}
 =\frac{2\alpha}{\Delta x^2}
 \left(T_{N-1}-T_N-\frac{\Delta x}{k}q_{\mathrm{out}}\right).
\end{equation}
The same expression follows from an energy balance on a half-width end
cell. Per unit cross-sectional area, its heat capacity is
$\rho c_p\Delta x/2$, and its net heat input is
$k(T_{N-1}-T_N)/\Delta x-q_{\mathrm{out}}$.
The factor of two reflects that smaller end-cell volume. The ghost node
is eliminated algebraically; no extra array element is required.


\subsection*{Two choices of time level}
For explicit FTCS, evaluate the parentheses in
Equation~\eqref{eq:surface} at time $n$. For Backward Euler, evaluate
them at $n+1$ and move unknown temperatures to the matrix side.
Insulation therefore does \emph{not} require overwriting the surface
temperature with the adjacent node's value: the surface retains heat
capacity and receives its own update.


\clearpage
\section{Recipe A: specified right heat flux or insulation}
\textbf{Inputs for either method.} Near the other physical inputs, add:
\begin{lstlisting}[label=lst:fluxinput]
q_flux = 1000.0  # W/m^2, positive OUT of the right end
\end{lstlisting}
Use \code{q_flux = 0.0} for insulation. For the modified right surface,
delete \code{T_history[:, -1] = T_right} in the initialization block.
Keep \code{T_history[0, :] = T_initial} and
\code{T_history[:, 0] = T_left}. The initial right temperature is now
\code{T_initial}; its later values will be calculated.


\subsection*{Explicit file: replace its entire time loop with this block}
From Equation~\eqref{eq:surface},
\[
 T_N^{n+1}=T_N^n+2\Fo
 \left(T_{N-1}^n-T_N^n-\frac{\Delta x}{k}q_{\mathrm{flux}}\right).
\]
\begin{lstlisting}[label=lst:explicitflux]
for n in range(number_of_steps):
    for i in range(1, number_of_nodes - 1):
        T_history[n + 1, i] = T_history[n, i] + Fo * (
            T_history[n, i - 1]
            - 2.0 * T_history[n, i]
            + T_history[n, i + 1]
        )
    # Inside the time loop, OUTSIDE the spatial loop:
    T_history[n + 1, -1] = T_history[n, -1] + 2.0 * Fo * (
        T_history[n, -2] - T_history[n, -1]
        - dx * q_flux / k
    )
\end{lstlisting}
Keep \code{solver_start} before this loop and \code{solver_seconds}
after it. The original $\Fo\le0.5$ check remains suitable for this
specified-flux update.


\subsection*{Implicit file: change the last matrix row and its right side}
The new surface equation is
\[
 -2\Fo T_{N-1}^{n+1}+(1+2\Fo)T_N^{n+1}
 =T_N^n-2\Fo\frac{\Delta x}{k}q_{\mathrm{flux}}.
\]
Replace \code{A[-1, -1] = 1.0} before the time loop with
\begin{lstlisting}[label=lst:implicitfluxmatrix]
A[-1, :] = 0.0  # Clear the entire old boundary row.
A[-1, -2] = -2.0 * Fo
A[-1, -1] = 1.0 + 2.0 * Fo
\end{lstlisting}
Inside the time loop, replace \code{b[-1] = T_right} with the following
line, indented four spaces like \code{b[0] = T_left}:
\begin{lstlisting}[label=lst:implicitfluxrhs]
    b[-1] = T_history[n, -1] - 2.0 * Fo * dx * q_flux / k
\end{lstlisting}
Keep the interior matrix rows and the left boundary row as written.
Inside the time loop, also keep \code{b = T_history[n, :].copy()} and the call to
\code{np.linalg.solve(A,b)}.


\clearpage
\section{Recipe B: convection at the right surface}
Start from a copy with the original fixed boundaries, not a mixture of
Recipes A and B. Add these inputs near the top:
\begin{lstlisting}[label=lst:robininput]
h = 100.0          # W/(m^2 K), convection coefficient
T_infinity = 20.0  # deg C, surrounding fluid temperature
\end{lstlisting}
After \code{dx} and \code{Fo} have been calculated, add
\begin{lstlisting}[label=lst:biot]
if h < 0:
    raise ValueError("h must be nonnegative.")
Bi = h * dx / k
\end{lstlisting}
Delete \code{T_history[:, -1] = T_right} from initialization, leaving
the right surface initially at \code{T_initial}.


\subsection*{Explicit update and its extra stability check}
Substitute $q_{\mathrm{out}}^n=h(T_N^n-T_\infty)$ into the flux update:
\[
 T_N^{n+1}=T_N^n+2\Fo
 \left[T_{N-1}^n-(1+\Bi)T_N^n+\Bi T_\infty\right].
\]
Use the full explicit loop from Recipe A, but replace its last surface
assignment with this block (including the shown indentation):
\begin{lstlisting}[label=lst:explicitrobin]
    T_history[n + 1, -1] = T_history[n, -1] + 2.0 * Fo * (
        T_history[n, -2] - (1.0 + Bi) * T_history[n, -1]
        + Bi * T_infinity
    )
\end{lstlisting}
Do not also use \code{q_flux} at this boundary. The coefficient of the
old surface temperature is $1-2\Fo(1+\Bi)$. A sufficient stability
restriction giving nonnegative update coefficients is
\[
 \Fo\le\frac{1}{2(1+\Bi)}.
\]
After defining \code{Bi}, add the following check before the time loop.
The existing interior check can stay in place.
\begin{lstlisting}[label=lst:robincheck]
Fo_limit = 1.0 / (2.0 * (1.0 + Bi))
if Fo > Fo_limit:
    raise ValueError("Reduce dt for the convection boundary.")
\end{lstlisting}
For the example $h=100$ and the current fine grid, $\Bi\approx0.001111$
and the limit is approximately $0.499445$, above the chosen $\Fo$.


\subsection*{Implicit matrix row and right-hand side}
Using $q_{\mathrm{out}}^{n+1}=h(T_N^{n+1}-T_\infty)$ gives
\[
 -2\Fo T_{N-1}^{n+1}+[1+2\Fo(1+\Bi)]T_N^{n+1}
 =T_N^n+2\Fo\Bi T_\infty.
\]
Replace the original \code{A[-1, -1] = 1.0} with
\begin{lstlisting}[label=lst:implicitrobinmatrix]
A[-1, :] = 0.0
A[-1, -2] = -2.0 * Fo
A[-1, -1] = 1.0 + 2.0 * Fo * (1.0 + Bi)
\end{lstlisting}
Replace \code{b[-1] = T_right} inside the time loop with
\begin{lstlisting}[label=lst:implicitrobinrhs]
    b[-1] = T_history[n, -1] + 2.0 * Fo * Bi * T_infinity
\end{lstlisting}
For constant $h$, the matrix is still built once. Backward Euler does not
need the explicit convection time-step restriction for this linear problem.


\section{What changes at the left end?}
The outward direction is $-x$ at the left end. Therefore, with outward
heat flux positive,
\[
 q_{\mathrm{out,left}}=+k\left.\frac{\partial T}{\partial x}\right|_0,
 \qquad +k\left.\frac{\partial T}{\partial x}\right|_0
 =h_{\mathrm{left}}(T_0-T_{\infty,\mathrm{left}}).
\]
Using this outward convention, the right-surface formulas apply at the
left with the surface index $N$ replaced by $0$ and its neighbor $N-1$
replaced by $1$. Positive outward flux always removes heat.


For example, to replace a fixed \emph{left} end with specified outward
flux, define \code{q_left} near the inputs, delete
\code{T_history[:, 0] = T_left}, and initialize the left surface with
\code{T_initial}. In the explicit time loop, outside its inner loop, use
\begin{lstlisting}
    T_history[n + 1, 0] = T_history[n, 0] + 2.0 * Fo * (
        T_history[n, 1] - T_history[n, 0] - dx * q_left / k
    )
\end{lstlisting}
In the implicit code, replace \code{A[0, 0] = 1.0} with
\begin{lstlisting}
A[0, :] = 0.0
A[0, 0] = 1.0 + 2.0 * Fo
A[0, 1] = -2.0 * Fo
\end{lstlisting}
and replace \code{b[0] = T_left} with
\begin{lstlisting}
    b[0] = T_history[n, 0] - 2.0 * Fo * dx * q_left / k
\end{lstlisting}
For left convection, define \code{Bi_left = h_left * dx / k}. The explicit
parentheses become
$T_1^n-(1+\Bi_{\mathrm{left}})T_0^n+
\Bi_{\mathrm{left}}T_{\infty,\mathrm{left}}$.
The implicit main-diagonal coefficient becomes
$1+2\Fo(1+\Bi_{\mathrm{left}})$ and
$b_0=T_0^n+2\Fo\Bi_{\mathrm{left}}T_{\infty,\mathrm{left}}$.
If both ends convect, apply the explicit restriction for both ends, using
the smaller allowable Fourier number.


\textbf{Do not silently change flux conventions.} If a problem instead
defines left flux by $-k\,\partial T/\partial x=q_x$, positive $q_x$
points into the rod. In that case $q_{\mathrm{out,left}}=-q_x$, so the
left flux contribution has a plus sign. State the convention before editing.


\section{Keep the contour colors consistent with the new problem}
After changing a boundary, \code{T_right} or \code{T_left} may no longer
control that surface. Flux or a different ambient temperature may also
produce temperatures outside the original $20$--$100\,^{\circ}\mathrm C$
range. In the plotting section, replace the two original min/max lines by
\begin{lstlisting}[label=lst:plotrange]
T_min = float(np.min(T_history))
T_max = float(np.max(T_history))
\end{lstlisting}
Keep the uniform-temperature fallback, \code{levels}, \code{contourf},
and PNG saving code. This covers all computed values. For side-by-side
method comparisons, choose identical color limits covering both results.
The plot's vertical axis is \emph{time}, not a second spatial direction.


\section{Checks that tell you whether an edit makes sense}
\subsection*{Check against a steady solution}
At steady state, the temperature is a straight line because
$d^2T/dx^2=0$. With left temperature fixed, the following targets are useful:
\begin{align*}
\text{Fixed right temperature:}\quad
 T(x)&=T_{\mathrm{left}}+
 (T_{\mathrm{right}}-T_{\mathrm{left}})\frac{x}{L},\\
\text{Specified right outward flux:}\quad
 T(x)&=T_{\mathrm{left}}-\frac{q_{\mathrm{flux}}}{k}x,\\
\text{Right convection:}\quad
 q_{\mathrm{steady}}&=
 \frac{T_{\mathrm{left}}-T_\infty}{L/k+1/h},\qquad h>0,\\
 T(x)&=T_{\mathrm{left}}-\frac{q_{\mathrm{steady}}}{k}x.
\end{align*}
For an insulated right end, the eventual temperature is uniformly
$T_{\mathrm{left}}$. For the example flux $1000\ \mathrm{W/m^2}$,
the steady right temperature is approximately $97.78\,^{\circ}\mathrm C$.
For $h=100\ \mathrm{W/(m^2\,K)}$ and $T_\infty=20\,^{\circ}\mathrm C$,
it is approximately $85.45\,^{\circ}\mathrm C$.
These are \emph{long-time} targets, not temperatures guaranteed at 300 s.
Increase \code{total_time} and check whether later profiles stop changing.


\subsection*{Simple experiments}
\begin{enumerate}
\item Set both fixed temperatures and the initial temperature equal.
      Both methods should retain a uniform temperature at every time.
\item Set the right flux to zero with the left end held hot. The right
      temperature should evolve; it should not stay fixed at its old value.
\item Set $h=0$ in the convection recipe. Its equations should become
      exactly the insulated-boundary equations.
\item Use a positive right outward flux. Reversing its sign should reverse
      the direction of that boundary's imposed heat transfer.
\item Reduce $\Delta x$ and $\Delta t$ consistently in both methods.
      Compare temperatures at the same physical position and time;
      their difference should decrease as the solution is resolved.
\end{enumerate}


\subsection*{Common beginner errors and their fixes}

\textbf{The new boundary remains fixed.}
Remove its whole-column assignment. In the implicit file, also replace its
matrix row and its \code{b} entry.

\textbf{IndentationError.}
Use spaces consistently. A boundary update belongs inside the time loop, at
the same indentation as \code{for i}, not inside the spatial loop. An implicit
\code{b} change belongs before the solve.

\textbf{NameError.}
Define new inputs before using them, and define \code{Bi} after \code{dx}.
Python runs the file from top to bottom.

\textbf{Temperatures oscillate or explode.}
Recheck the explicit time-step restriction, including convection, and make
sure all right-side values use the old time row.

\textbf{The image has blank temperature bands.}
Use the updated plotting limits from the preceding section so the levels cover
the computed range.

Save the file after each edit. Read the last line of a Python error message
first. It often identifies the problem and points to the relevant line.
\end{document}


